% Part 5: the black-hole information problem read in IAM. Carries the information-paradox paper (Feb 2026; 453 PDF lines, read in full
% 2026-10-02) Sections 1, 2, 4, 5, 8.2-8.3, 9 and 10, and Section 5 of the black-hole thermodynamics paper (25 Feb 2026; 431 lines).
% Its thermodynamic Sections 3, 6 and 7 are carried in Chapter ch:blackholes; they are not repeated here.
% Corrections applied: PAPER_ERRATA B1, B2, B6 (fixed point not carried), B7 (island row; drift records M(t) only), B8 (no claim to resolve),
% M1 (M-sigma removed), L1 (superposition/record wording); BLACK_HOLES_CHECK #3, #4, #8, #9, #10.
% Every claim of the projection picture is labelled \conjecture or \interp: it is the author's reading, not a derivation.

\chapter{The black-hole information problem read as a record on the horizon}\label{ch:bhinformation}

\section*{Overview}
This chapter sets out a conceptual reframing of the black-hole information problem within IAM, the gravitational-decoherence framework whose
cosmological sector is tested against the Planck 2018 CMB likelihood in the chains of Chapter~\ref{ch:dual}. The thermodynamic foundation --- the black-body identity
$P_{\rm SB}=P_{\rm Hawking}$, the encoding rate $\Gamma_{\rm BH}=c^3/(1920\,GM\ln2)$, and the two-horizon framework --- is established in
Chapter~\ref{ch:blackholes}; this chapter draws on those results.

The central idea is a physical picture: matter does not enter a black hole --- the projection ceases. In IAM, a particle exists as a stable decoherence loop
registered on the cosmic holographic boundary. The black-hole horizon is the surface at which that registration becomes causally impossible. Crossing the
horizon does not mean travelling inward; it means the decoherence loop that constitutes the particle can no longer be maintained. What remains on the
horizon is the entropy imprint --- the pixel structure of the last recorded state --- not an interior object carrying information. \conjecture

The picture addresses the information paradox by questioning its premise: on this reading, information was never interior. No exotic machinery is
introduced and no firewall arises; the record is encoded on the horizon surface at the moment of crossing, consistent with Hawking's final ``soft hair''
position~\cite{HawkingPerryStrominger2016}, for which IAM offers a thermodynamic mechanism. Computing the fine-grained entropy of the radiation, and carrying the first-law entropy transfer of Chapter~\ref{ch:blackholes} to the Page curve, are the open steps. What the picture offers is a
reading in which the question changes. \interp

\section{Introduction}\label{sec:bi_intro}
The black-hole information paradox has resisted resolution for fifty years despite contributions from Hawking, Susskind, Maldacena, Penrose, Polchinski and
many others~\cite{Hawking1976,Page1993,Maldacena1998,AMPS2013,MaldacenaSusskind2013,Penrose1989}. It is worth noting that Hawking himself substantially revised his position over the decades. His
original argument~\cite{Hawking1975,Hawking1976} held that thermal Hawking radiation destroys information, violating unitarity. By 2004 he conceded that information is likely
preserved~\cite{Hawking2005}. In his final work --- the ``soft hair'' papers with Perry and Strominger~\cite{HawkingPerryStrominger2016} --- he argued that information is stored on the
horizon surface itself, not in the interior, encoded in zero-energy quantum excitations he called soft hair. IAM does not contradict this final position. It
offers the thermodynamic mechanism that Hawking's work was pointing toward: an account of why information must be on the surface, grounded in the Landauer
principle and gravitational decoherence. \interp

IAM starts from the observation that gravitational decoherence --- the irreversible conversion of quantum superpositions to classical outcomes --- produces
information at a thermodynamic cost governed by the Landauer principle~\cite{Landauer1961} (Chapter~\ref{ch:iams_law}). This cost is paid at the cosmic Gibbons--Hawking
temperature $T_{\rm GH}$ and encoded on the cosmic holographic boundary. The same mechanism, operating locally rather than globally, produces black holes when
local decoherence exceeds the encoding capacity of the local bounding surface (Chapter~\ref{ch:blackholes}, Section~\ref{sec:bh_formation}).

The chapter establishes one central physical picture and draws its consequences. Section~\ref{sec:bi_entropy} introduces the conceptual foundation: the distinction between
entropy (pixel capacity) and information (pixel content), and the picture in which matter does not enter a black hole but ceases to be projected at the
horizon. Section~\ref{sec:bi_surfaces} presents the reframing of black holes as encoding surfaces. Section~\ref{sec:bi_paradox} reads the information paradox from its premise,
and Section~\ref{sec:bi_twohorizon} gives the same reading in the two-horizon framework. Section~\ref{sec:bi_astro} collects two astrophysical consequences, and
Section~\ref{sec:bi_compare} compares the picture with existing approaches.

\section{Entropy, information, and the cessation of projection}\label{sec:bi_entropy}
Before the consequences, it is worth establishing a conceptual foundation that shapes how all subsequent statements should be read. Two distinctions matter
here, and conflating them has contributed to fifty years of confusion about what the black-hole information paradox is actually asking.

\subsection{Entropy is not information}
In the context of black holes, these two words are often used interchangeably. They are not the same thing.

\emph{Entropy} ($S_{\rm BH}=\kB A/4\lP^2$) is a count of available encoding slots --- the number of pixels on the holographic surface, the structural capacity of the
boundary. It tells you how many distinct states could be encoded. It says nothing about what is actually written.

\emph{Information} is what is recorded in those slots --- the specific content of the final state, the record of which quantum superpositions were resolved and
how. In IAM language: entropy is the pixel array; information is what is written on the pixels.

Decoherence is the act of writing: it converts a quantum superposition into a classical record at the Landauer cost of $\kB T\ln2$ per bit. The Landauer
principle bridges entropy and information exactly --- each bit written costs $\kB T\ln2$ of energy at the local temperature.

This distinction matters because it changes what we think is stored on the black-hole horizon. The Bekenstein--Hawking entropy $S_{\rm BH}=\kB A/4\lP^2$ is the
capacity of the horizon --- the number of pixels available, one nat per $4\lP^2$ and one bit per $4\ln2\,\lP^2$ (Chapter~\ref{ch:bekenstein}). The information encoded in
those pixels is the specific recorded state of the last decoherence event before the projection was severed. The horizon does not store information in the
way a hard drive stores a file. It stores the entropy structure --- the pixel pattern --- that is that information, frozen at the moment projection ceased.
\interp

\subsection{Matter does not enter a black hole --- the projection ceases}
This is the central physical claim of the picture, and it addresses the information paradox more directly than any mechanism that accepts the premise of
interior storage.

In IAM, a particle exists as a stable, self-consistent decoherence loop. Three elements are required to maintain it:
\begin{enumerate}
\item \emph{The superposition} --- the quantum state of the field, not yet recorded.
\item \emph{The projector} --- gravitational decoherence, the process that converts the superposition to a classical outcome at the Landauer cost.
\item \emph{The screen} --- the cosmic holographic boundary, the encoding surface on which the recorded state is registered. Without registration, the
decoherence loop cannot be maintained.
\end{enumerate}
A particle exists because its decoherence loop is continuously refreshed and registered on the cosmic encoding surface. The Landauer cost is the energy
required to keep the projector running. On this reading the rest mass of the particle is this cost, evaluated at the Gibbons--Hawking temperature; the
electron-mass relation of Chapter~\ref{ch:electronmass} is the quantitative form of that statement for the electron. \conjecture

When a particle approaches a black-hole horizon, something specific happens: the cosmic encoding surface becomes causally inaccessible. The interior of a
black hole is not connected to the cosmic horizon by any causal path. The screen goes dark.

The projector cannot reach the screen. The decoherence loop cannot be registered. The particle ceases to be projected.

It does not travel inward and carry its information with it. It does not get copied onto the horizon by some exotic mechanism. The loop that constituted
it simply cannot be sustained in a region where the cosmic encoding surface is inaccessible.

What appears on the black-hole horizon is not the particle. It is the entropy imprint of the last recorded state --- the final frame of the projection before
the screen went dark. This entropy imprint, frozen in the pixel structure of the horizon at $S_{\rm BH}=\kB A/4\lP^2$, encodes the informational content of that
final state. The Bekenstein--Hawking entropy is not a measure of hidden interior degrees of freedom. It is the area of the last frame. \conjecture

This picture is consistent with Hawking's final ``soft hair'' result~\cite{HawkingPerryStrominger2016} and offers a mechanism for it. Soft hair is, on this reading, the pixel
structure of the last frame, written in zero-energy quantum excitations on the horizon surface. IAM explains why that structure must be there: the Landauer
principle requires that every decoherence event be registered on an encoding surface, and the horizon is the last surface the projector could reach.
\conjecture

The information paradox, in this picture, is not a paradox about escaping information. It is a question about what it means for matter to ``enter'' a black
hole at all. The answer IAM offers is: it does not. The projection stops. The entropy imprint remains. That imprint is slowly radiated away as Hawking
radiation as the horizon shrinks --- not because information escapes from the interior, but because the last frame is being erased one pixel at a time, each
erasure costing $\kB T_{\rm BH}\ln2$ of energy, exactly as Landauer requires. Summed over the horizon, those erasures cost $T_{\rm BH}S_{\rm BH}=\tfrac12Mc^2$, the Smarr
half of Chapter~\ref{ch:blackholes}. \interp

\section{Black holes as encoding surfaces}\label{sec:bi_surfaces}
\subsection{The reframing}
\emph{Standard picture:} gravity collapses matter; extreme curvature creates a one-way membrane; nothing escapes.

\emph{IAM picture:} gravitational decoherence accelerates during collapse; the information production rate exceeds local encoding capacity; the universe
resolves the encoding deficit by creating a new surface. The horizon is not a barrier --- it is an encoding surface created by thermodynamic necessity.
\conjecture

This reframing has a specific consequence: the black-hole horizon is not the boundary between ``inside'' and ``outside.'' It is the storage medium. The
information is not behind the horizon --- it is on the horizon.

\subsection{The saturation condition}
The Bekenstein--Hawking entropy $S_{\rm BH}=\kB A/(4\lP^2)$ is a surface quantity, not a volume quantity. The maximum information storable in any region is bounded
by the area of its boundary (the holographic principle~\cite{tHooft1993,Susskind1995}), not its volume. On the IAM reading, the interior of a black hole has no
independent degrees of freedom beyond those encoded on the horizon.

This is the IAM answer to the information paradox stated as a premise correction rather than a mechanism: information does not enter the interior and
require extraction. Information is written onto the horizon surface as matter crosses it, at the Landauer cost of $\kB T_{\rm BH}\ln2$ per bit. There is no
interior information to recover. \conjecture\ The formation criterion itself, saturation of the bound and its equality with the hoop conjecture, is
Chapter~\ref{ch:blackholes}, Eq.~\eqref{eq:bh_hoop}; the two regimes of the mechanism, self-regulating and runaway, are the table of
Section~\ref{sec:bh_regimes}.

\section{The information paradox read from its premise}\label{sec:bi_paradox}
\subsection{The paradox and its premise}
Hawking~\cite{Hawking1975} showed that black holes emit thermal radiation. Thermal radiation is random: it carries maximum entropy and no information about what produced
it. If a black hole fully evaporates via thermal Hawking radiation, the information about its formation is gone. This violates quantum-mechanical
unitarity~\cite{Hawking1976}.

Every proposed resolution --- complemen\-tarity~\cite{SusskindThorlaciusUglum1993}, AdS/CFT~\cite{Maldacena1998}, firewalls~\cite{AMPS2013}, ER\,=\,EPR~\cite{MaldacenaSusskind2013}, the island
formula~\cite{Penington2020,Almheiri2020} --- accepts that the information must be recovered from somewhere, and works out machinery to extract it.

\subsection{The premise in question}
\begin{keybox}[title={Proposition (no interior degrees of freedom) --- conjecture}]
The interior of a black hole has no degrees of freedom independent of those encoded on the horizon surface.
\end{keybox}
\emph{Argument.} The Bekenstein--Hawking entropy $S_{\rm BH}=\kB A/(4\lP^2)$ scales with area, not volume. The holographic principle~\cite{tHooft1993} states that the maximum
information content of any region is bounded by its boundary area. For a black hole, the boundary is the horizon. The interior is then the holographic
projection of the boundary state, not an independent information storage system, and there is no interior information to be lost. \conjecture\ The
argument promotes a bound on capacity to a statement about where degrees of freedom live; that step is the content of the holographic principle, which
is established in anti-de Sitter space~\cite{Maldacena1998} and assumed here for the black holes of our universe.

\subsection{Evaporation as the surface shrinking}
If information is on the surface, then evaporation is the surface shrinking. As the black hole radiates:
\begin{itemize}
\item $M$ decreases; $A_{\rm BH}$ decreases; $S_{\rm BH}$ decreases.
\item Each bit of surface entropy is paid off in the Hawking luminosity at the Landauer cost $\kB T_{\rm BH}\ln2$ (Chapter~\ref{ch:blackholes}, Section~\ref{sec:bh_bitrate}).
\item The information is not ``escaping'' the interior --- the surface is transferring it to the outgoing radiation.
\end{itemize}
The radiation is thermal at each instant, but the temperature $T_{\rm BH}(M)$ rises as $M$ decreases, so the time series of Hawking emissions carries a systematic
frequency shift. That shift records the evaporation history $M(t)$, which depends on the mass alone; it does not by itself carry the information about what
fell in. If the information is in the radiation, it must be in correlations across time that a thermal spectrum at each instant does not show. \derived\ (the
drift) \conjecture\ (correlations)

On this reading no information is destroyed, no unitarity is violated, no firewall is required and no non-locality is needed: the paradox dissolves because
its premise --- that information is interior --- is not correct. \interp\ The next step is the calculation that would show it: the fine-grained entropy of the radiation, which for a unitary evaporation must rise and then fall to zero (the Page curve~\cite{Page1993}). The first-law entropy transfer
of Chapter~\ref{ch:blackholes}, Eq.~\eqref{eq:bh_Str_closed}, rises monotonically to $S_{\rm BH,0}$ and reaches half of it at $0.646\,\tau_{\rm evap}$; it is a bound on the Page curve,
not the Page curve itself (Fig.~\ref{fig:bh_transfer}). \openprob

\subsection{Why the firewall does not arise}
The AMPS argument~\cite{AMPS2013} shows that if late Hawking radiation carries information, quantum monogamy is violated: the outgoing radiation cannot be maximally
entangled with both the early radiation (required for unitarity) and the infalling modes (required by the equivalence principle).

The IAM picture eliminates one of the three parties. The infalling modes are not a separate system --- they are encoded on the horizon surface at the moment
of crossing. There is no independent interior mode to participate in the monogamy paradox. The AMPS trilemma then has only two parties, which is not a
paradox. \conjecture

\section{The same reading in the two-horizon framework}\label{sec:bi_twohorizon}
Hawking's original calculation~\cite{Hawking1975} suggested that information falling into a black hole is irretrievably lost upon evaporation, violating unitarity. The paradox
arises from treating the black-hole horizon as a one-way membrane that destroys information.

In the two-horizon framework of Chapter~\ref{ch:blackholes} (Section~\ref{sec:bh_twohorizons}) the paradox is read on three grounds. First, information is not destroyed --- it is
transferred from a hot local encoding surface (the black-hole horizon) to a cooler global encoding surface (the cosmic horizon), exactly as heat transfers from a
hot body to a cold reservoir. Second, the transfer obeys the second law: energy and information flow from hot to cold at the thermally limited rate
$\Gamma_{\rm thermal}$. Third, the cosmic horizon is the final repository --- all locally encoded information migrates to the global ledger as the universe approaches
thermodynamic equilibrium. \interp

The reading is not that information escapes through subtle correlations in Hawking radiation~\cite{Page1993}. It is that information was never trapped: it sits on a
thermal surface that radiates by ordinary thermodynamics, transferring its contents to a cooler surface at the rate set by the temperature difference.
\interp\ As stated in Chapter~\ref{ch:blackholes} (Section~\ref{sec:bh_scope}), the second law fixes the
direction and rate of the coarse-grained flow; whether the fine-grained state is preserved is an open question for the framework.

\section{Astrophysical consequences}\label{sec:bi_astro}
\subsection{Early black holes}
JWST has identified massive black holes at $z>7$, earlier than standard accretion models easily explain. In IAM, black holes are the universe's primary
encoding infrastructure before large-scale structure formation makes the cosmic horizon dominant (Chapter~\ref{ch:blackholes}, Section~\ref{sec:bh_formation}; Chapter~\ref{ch:theoryinterp}).
Early, massive black-hole formation is then not a puzzle but an expectation: the universe front-loads its encoding capacity because black holes are the most
efficient encoding surfaces available in the early universe. \conjecture

The IAM expectation is that early black-hole masses should correlate with the local information-production environment (density, interaction rate) rather
than with accretion history alone. \conjecture\ It becomes a prediction when the collapse entropy of Table~\ref{tab:bh_seed} is computed for a given halo. \openprob

\subsection{Cosmic censorship as a Landauer requirement}
Penrose's cosmic censorship conjecture~\cite{Penrose1969} --- that naked singularities cannot form in generic collapse --- is an independent postulate in general relativity.
In IAM it reads as a consequence of the Landauer principle: a naked singularity would be a region of unbounded information production with no encoding
surface. Landauer's principle is not optional: information must be encoded. The universe must provide a horizon. On this reading cosmic censorship is
thermodynamically enforced, not geometrically postulated. \conjecture

\section{Comparison with existing approaches}\label{sec:bi_compare}
\begin{table}[htbp]\centering\small
\caption{The IAM reading beside existing approaches to the information problem. \interp}\label{tab:bi_compare}
\begin{tabularx}{\textwidth}{@{}>{\raggedright\arraybackslash}p{0.25\textwidth}>{\raggedright\arraybackslash}X>{\raggedright\arraybackslash}X@{}}\toprule
approach & core mechanism & IAM reading\\\midrule
complementarity~\cite{SusskindThorlaciusUglum1993} & observer-dependent description & simpler: the information is on the surface for all observers\\
AdS/CFT~\cite{Maldacena1998} & the boundary encodes the bulk (anti-de Sitter) & the same conclusion, applied to the de Sitter universe we live in, without a dual theory\\
firewall~\cite{AMPS2013} & drama at the horizon preserves unitarity & no independent interior degrees of freedom; the firewall is unnecessary\\
ER\,=\,EPR~\cite{MaldacenaSusskind2013} & Hawking pairs linked by wormholes & no wormholes; ordinary thermodynamics\\
island formula~\cite{Penington2020,Almheiri2020} & replica wormholes give a Page curve that turns over and falls to zero & no topology; the first-law transfer gives only the monotonic coarse-grained curve and its bound (Fig.~\ref{fig:bh_transfer}), not the turnover\\\bottomrule
\end{tabularx}
\end{table}
The IAM reading is the most economical of these: it requires no new mathematical structures, no exotic topology, no modifications to quantum mechanics, and
no non-locality. It requires only that the holographic principle be taken seriously --- that the horizon surface is the storage medium, not a boundary
concealing interior storage. The island formula computes the fine-grained entropy and obtains the turnover; the next step for the IAM
reading is to obtain it as well. \interp

\section{Discussion}\label{sec:bi_discussion}
Three statements have been set out. (1) In the black-body approximation the Stefan--Boltzmann luminosity of the horizon and the Hawking luminosity are one
formula, and the horizon writes and erases bits at $\Gamma=c^3/(1920\,GM\ln2)$ (Chapter~\ref{ch:blackholes}). \derived\ (2) The information paradox changes character when the
premise that information is interior is replaced by the holographic statement that information is on the surface. \interp\ (3) The first-law entropy
transfer during evaporation follows analytically from the Landauer emission rate with no free parameters; it bounds the Page curve, reaching half the
initial horizon entropy at $0.646\,\tau_{\rm evap}$, but is not the Page curve. \derived

The structural link is the universal temperature formula $T=\hbar\kappa/(2\pi\kB c)$, the same $2\pi$ at the black-hole, cosmic and Rindler horizons
(Chapter~\ref{ch:blackholes}, Section~\ref{sec:bh_thermo}; Chapter~\ref{ch:bekenstein}). It is what lets one accounting run from the horizon of a black hole to the horizon of the universe.
IAM reads cosmology (Part~\ref{part:2}), the particle scale (Chapters~\ref{ch:koide}, \ref{ch:electronmass} and~\ref{ch:higgsrecord}) and black holes with the same mechanism, seen from different angles. \interp

Two open directions follow from that link. First, the encoding throughput $\Gamma_{\rm BH}=c^3/(1920\,GM\ln2)$, a Landauer rate at a horizon, is the natural template
for deriving the $\alpha^{5/2}$ suppression in the electron-mass relation of Chapter~\ref{ch:electronmass}, where that suppression is still an open problem. Second, the $2\pi$
in $T_{\rm BH}$ is the same Euclidean $2\pi$ as in $T_{\rm GH}$, and the prefactor $(2\pi)^{-1/10}$ of the electron relation, whose factor $(2\pi)^{3/10}$ beyond the $T_{\rm GH}$ part was found by numerical search, may be that same
$2\pi$. Completing either derivation would deepen the black-hole thermodynamic framework and the particle-scale one at once. \conjecture

One near-term consequence of the picture concerns the two sectors. In the dual-sector picture photons ($\Sigma=1$) do not couple to the record term and
matter does. Decoherence of matter superpositions should therefore carry a component that scales with the local gravitational environment, and photonic
superpositions should not. \conjecture\ For present quantum hardware the gravitational times are far beyond reach under any model (Chapter~\ref{ch:gravdec},
Section~\ref{sec:gd_sectors}), so the test needs massive superpositions.

This reading turns the information question into a question about records and the surfaces that hold them. A quantum-device physicist working with massive superpositions is aimed at exactly the regime where it can be tested.
